\documentclass{article}
\usepackage[a4paper, total={6.5in, 10in}]{geometry}
\setlength{\parindent}{0pt}
\usepackage{tcolorbox}
\tcbset{colback=white!94!black, colframe=black!60!white, coltitle = white, boxrule=0.8pt, arc=2mm}
\newtcolorbox{boxs}[2][]{title= #2,#1}
\usepackage{datetime2}
\usepackage[british]{babel}
\usepackage{amsfonts}
\usepackage{amsmath}

\title{\textbf{Programming HW 2}}
\author{Robert O'Connell}
\date{\today}
\begin{document}
\maketitle

\section*{Question 1}
\begin{center}
\begin{tabular}{|c|c|c|c|c|c|c|c|}
\hline
$A$ & $B$ & $C$ & $D$ & $A \ \text{AND} \ B$ & $\text{NOT} \ C$ & $\text{NOT} \ C \ \text{AND} \ D$ & $(A \ \text{AND} \ B)\ \  \text{OR}\ \ (\text{NOT} \ C \ \text{AND} \ D)$ \\ \hline
F & F & F & F & F  & T  & F  &  F  \\ \hline
F & F & F & T & F  & T  & T  & T  \\ \hline
F & T & F & F &  F & T  & F  &  F \\ \hline
F & T & F & T &  F & T  & T  &  T \\ \hline
T & F & F & F & F  & T  & F  &  F \\ \hline
T & F & F & T &  F & T  & T  &  T \\ \hline
T & T & F & F &  T & T  & F  & T  \\ \hline
T & T & F & T &  T & T  & T  & T  \\ \hline
F & F & T & F &  F & F  & F  &  F \\ \hline
F & F & T & T &  F & F  & F  & F  \\ \hline
F & T & T & F &  F & F  & F  & F  \\ \hline
F & T & T & T &  F & F  & F  & F  \\ \hline
T & F & T & F & F  & F  & F  & F  \\ \hline
T & F & T & T & F  &  F & F  & F  \\ \hline
T & T & T & F &  T & F  &  F & T  \\ \hline
T & T & T & T &  T & F  & F  & T  \\ \hline
\end{tabular}
\end{center}


\section*{Question 2}
\[
(n)(n+1)(n+2)(n+3) = 24k, \quad k \in \mathbb{Z}, n \in \mathbb{Z}, 
\]
\subsubsection*{Proof}
\[
P = (n)(n+1)(n+2)(n+3)
\]
Since we are looking at the products of four consecutive integers we know that two are divisible by two, one of those two also divisible by four 
and another by three

Thus \(P\) is divisible by 2,3 and 4.
If \(P\) is divisible by 2,3 and 4 it must also be divisible by \(2\times 3 \times 4 = 24\)


\section*{Question 3}
\subsubsection*{Proof by contrapositive}
Trying to prove \(p \rightarrow q\)

Contrapositive is NOT \(q \rightarrow\) NOT \(p \)

\subsubsection*{Restatement of Question}
\(2^{n}-1\) is prime \(\rightarrow \ n\) is prime

\(n\) is NOT prime \(\rightarrow \ 2^{n}-1\) is NOT prime

\subsubsection*{Proof}
\(n \in \mathbb{N}\)

Assume \(n\) is not prime and therefore can be written as a product of two integers, where \(1<a,b<n\)
\[
n = ab
\]
\[
2^{n}-1 = 2^{ab}-1
\]

Then using the expression \(x^{k}-1 = (x-1)((x)^{k-1}+(x)^{k-2} \dots x+1)\), we get
\[
    2^{n}-1 = (2^a-1)((2^a)^{b-1}+(2^a)^{b-2} \dots 2^a+1)
\]
Since \(a,b > 0 \), we have that \(2^n-1\) can be written as the product of two integers, meaning that if \(n\) is not a prime
then \(2^n-1\) is also not a prime.

Then the original statement, if \(2^n-1\) is prime then \(n\) is prime must be true


\section*{Question 4}
For any integers \(m\) and \(n\) where \(n \neq 0\), \(m+\sqrt{2}n\) is not rational

\subsubsection*{Step 1}

Assume \(n\sqrt{2}\) is rational,

\(n\sqrt{2} = \frac{a}{b}\) where the fraction is in its simplist form, i.e. highest common factor of \(a\) and \(b\) is 1 and \(b \neq 0\), \(a,b \in \mathbb{Z}\)

\[
2n^2 = \frac{a^2}{b^2}
\]
\[
2(b^2n^2) = a^2
\]
Thus \(a^2\) and \(a\) are even
\[
a = 2k, \quad z \in \mathbb{Z}
\]

\[
2(b^2n^2) = 4k^2
\]
\[
n^2b^2 = 2k^2
\]
Thus \(n^2b^2\) is even, and at least one of the factors must be even
\\

If \(n^2\) is odd then \(b^2\) and \(b\) must be even, which contradicts the assumption that the highest common 
factor of \(a,b\) is 1, a contradiction
\\

If \(n^2\) is even then \(n\) is even, then looking back at the original equation \(n\sqrt{2} = \frac{a}{b}\) we see that 
\(n\) and \(a\) are even, thus
\[
p\sqrt{2} = \frac{k}{b}
\]
This is a smaller fraction which represents the same number, a contradiction

Thus \(n\sqrt{2}\) must be irrational
\\

\subsubsection*{Step 2}
We now will check the case \(m+\sqrt{2}n\), knowing \(n\sqrt{2}\) is irrational

Assume \(m+\sqrt{2}n\) is rational, 

\( m + n\sqrt{2} = \frac{a}{b}\) where the fraction is in its simplist form, i.e. highest common factor of \(a\) and \(b\) is 1 and \(b \neq 0\), \(a,b \in \mathbb{Z}\)

We can reorder the equation to get
\[
n\sqrt{2} = \frac{a}{b} - m
\]
\[
n\sqrt{2} = \frac{a-mb}{b}
\]

This shows that \(n\sqrt{2}\) can be written as a fraction, which we know is false since we just proved it's irrationality

This is a contradiction, thus \(m + n\sqrt{2}\) is also irrational


\section*{Question 5}
\[
\sum_{i=0}^{n} i^2 = \frac{1}{6}n(n+1)(2n+1)
\]

\subsubsection*{Proof by Induction}
Check base case (\(n=0\))
\[
0 = \frac{1}{6}(0)(1)(2) = 0
\]

Assume formula holds for \(n = k\)
\[
\sum_{i=0}^{k} i^2 = \frac{1}{6}k(k+1)(2k+1)
\]


\begin{align*}
    \sum_{i=0}^{k} i^2 + (k+1)^2 &= \frac{1}{6}k(k+1)(2k+1) + (k+1)^2 \\
    \sum_{i=0}^{k+1} i^2 &= \frac{1}{6}(2k^3 + 3k^2 + k  + 6k^2 + 12k + 6 ) \\ 
    &= \frac{1}{6}(2k^3 + 9k^2 + 13k + 6 ) \\ 
    &= \frac{1}{6}(k+1)(2k^2 + 7k + 6) \\ 
    & = \frac{1}{6}(k+1)(k+2)(2k + 3) \\
    & = \frac{1}{6}(k+1)((k+1)+1)(2(k+1)+1)
\end{align*}

Thus the formula holds in the induction step, and the result is proven
\end{document}